\documentclass[12pt]{article}
\usepackage[utf8]{inputenc}
\usepackage[T1]{fontenc}
\usepackage[margin=1in]{geometry}
\usepackage{setspace}
\usepackage{xcolor}
\usepackage[colorlinks=true, urlcolor=blue!60!black, linkcolor=blue!60!black]{hyperref}
\setlength{\parindent}{2em}
\setlength{\parskip}{1.25em}
\renewcommand{\baselinestretch}{1.0}
\begin{document}
\begin{center}
{\LARGE \textit{Brown v.\ Board of Education (Brown I)}, 347 U.S.\ 483 (1954)\par}
\vspace{0.8em}
{\large Case Brief --- Four-Page Edition\par}
\end{center}
\vspace{1em}

\section{Memory Jogger}

A unanimous Court, per Chief Justice Warren, repudiated ``separate but equal'' in public education: state-sanctioned racial segregation of public schools is inherently unequal and violates the Equal Protection Clause (\href{https://supreme.justia.com/cases/federal/us/347/483/}{Justia}).

\section{Detailed Case Facts}

\textit{Brown} consolidated four class actions by Black schoolchildren from Kansas, South Carolina, Virginia, and Delaware challenging laws that required or permitted segregated public schools. In Topeka, Kansas, where state law permitted separate elementary schools, 13 parents (including Oliver Brown for his daughter Linda) sued after white schools refused their children; because the district court found the Black and white schools substantially equal, \textit{Brown} squarely challenged segregation itself. \textit{Briggs v.\ Elliott} (South Carolina) featured Dr.\ Kenneth Clark's psychological testimony; \textit{Davis} (Virginia) showed stark disparities; \textit{Gebhart v.\ Belton} (Delaware) challenged mandated segregation. A companion D.C.\ case, \textit{Bolling v.\ Sharpe}, was decided the same day under the Fifth Amendment.

The cases capped the NAACP Legal Defense Fund's campaign, led by Thurgood Marshall, against \textit{Plessy v.\ Ferguson} (1896), after graduate-school victories in \textit{Gaines}, \textit{Sipuel}, \textit{Sweatt}, and \textit{McLaurin}. The Justice Department urged invalidation as \textit{amicus}. After the 1952 argument before a divided Vinson Court, the Court ordered reargument on the Fourteenth Amendment's history; Vinson died, and Chief Justice Warren forged a unanimous Court.

\section{Procedural History}

Three-judge district courts in Kansas, South Carolina, and Virginia upheld segregation under \textit{Plessy} (Kansas while finding segregation harmful to Black children; South Carolina and Virginia while ordering equalization). In Delaware, the state courts ordered immediate admission because of tangible inequality, and state officials sought certiorari. The cases were argued in December 1952, reargued in December 1953, and decided May 17, 1954. The Court restored the cases to the docket for argument on remedies, resolved in \textit{Brown II} (1955).

\section{Judicial Votes}

\textbf{9--0}.
\begin{itemize}
\item \textbf{Majority:} Chief Justice Warren, joined by Justices Black, Reed, Frankfurter, Douglas, Jackson, Burton, Clark, and Minton.
\item \textbf{Concurrences and dissents:} None.
\end{itemize}

\section{Holding}

State laws requiring or permitting racial segregation in public schools violate the Equal Protection Clause. In public education, ``separate but equal'' has no place: segregating children solely by race deprives minority children of equal educational opportunities even when physical facilities and other tangible factors are equal. The holding is confined to \textit{de jure} segregation in public education; the Court did not reach the Due Process Clause and did not formally overrule \textit{Plessy} beyond schools, though later \textit{per curiam} rulings extended \textit{Brown} to beaches, buses, golf courses, and parks. Footnote 11's social science corroborated the Kansas finding rather than serving as a prerequisite for future claims.

\section{Analysis of Opinions}

\textbf{Unanimous Court (Warren).} Warren wrote a short, accessible, non-accusatory opinion. The history of 1868 was ``[a]t best, \ldots inconclusive'': public education was rudimentary then and nearly nonexistent for Black Americans. So the Court ``cannot turn the clock back to 1868'' and must assess education in light of its present place in American life. Education is the foundation of citizenship, and where the state provides it, it must be available to all on equal terms. Because tangible factors were equal or being equalized, the question was whether segregation \emph{itself} denies equal opportunity. Building on \textit{Sweatt} and \textit{McLaurin}, Warren held that intangible harms apply ``with added force to children in grade and high schools,'' and adopted the Kansas court's finding that segregation denotes inferiority, which Footnote 11 called ``amply supported by modern authority,'' citing Clark's doll studies and Myrdal's \textit{An American Dilemma}. Contrary language in \textit{Plessy} was rejected.

\textbf{Unanimity.} There were no separate opinions. Reed and Jackson had considered writing separately, and Frankfurter urged reargument partly to build consensus; the Justices saw that a fractured Court would invite defiance across the South.

\section{Examples: Future Applications}

\begin{enumerate}
\item \textbf{Same-Side: Segregated Parks, Beaches, or Buses.} A city requires separate but physically identical facilities. Courts invalidate the ordinance as a badge of racial inferiority, as the post-\textit{Brown} \textit{per curiam} orders did.
\item \textbf{Same-Side: Segregated State Universities.} A state runs a white flagship and a comparable Black university. Tangible parity does not cure state-enforced separation.
\item \textbf{Opposite-Side: \textit{De Facto} Imbalance.} Schools are racially imbalanced only because of private housing choices and race-blind zones. Without discriminatory state intent, there is no violation.
\item \textbf{Opposite-Side: Private Religious School.} A school that is not a state actor is not bound by the Equal Protection Clause, though statutes or tax rules (as in \textit{Bob Jones}) may still restrict it.
\item \textbf{Fence-Sitter: Voluntary Integration Plans.} A district uses race as a tiebreaker to prevent racial isolation. The plurality in \textit{Parents Involved} read \textit{Brown} as forbidding racial classification; Justice Breyer's dissent read it as targeting caste-enforcing exclusion.
\end{enumerate}

\section{Critique}

\textbf{Neutral principles.} Herbert Wechsler argued that the Court failed to articulate a neutral principle for preferring Black citizens' freedom to associate over white citizens' freedom not to; Charles Black replied that segregation was a system designed to stamp Black Americans as an inferior caste. \textbf{Interest convergence and remedy.} Derrick Bell argued that unanimity reflected a convergence with Cold War elite interests; others faulted the reliance on stigma findings and \textit{Brown II}'s ``with all deliberate speed,'' which enabled Massive Resistance.

\textbf{Originalism.} Raoul Berger read Warren's refusal to ``turn the clock back to 1868'' as a departure from original intent. Michael McConnell, Robert Bork, and Steven Calabresi defend the result on originalist grounds, while Michael Klarman doubts those reconstructions.

\section{Key Quotations}

\begin{quote}
``Today, education is perhaps the most important function of state and local governments. Compulsory school attendance laws and the great expenditures for education both demonstrate our recognition of the importance of education to our democratic society. It is required in the performance of our most basic public responsibilities, even service in the armed forces. It is the very foundation of good citizenship. Today it is a principal instrument in awakening the child to cultural values, in preparing him for later professional training, and in helping him to adjust normally to his environment. In these days, it is doubtful that any child may reasonably be expected to succeed in life if he is denied the opportunity of an education. Such an opportunity, where the state has undertaken to provide it, is a right which must be made available to all on equal terms.''
\end{quote}
\noindent --- Chief Justice Warren, for a unanimous Court, 347 U.S.\ at 493.

\begin{quote}
``We conclude that in the field of public education the doctrine of `separate but equal' has no place. Separate educational facilities are inherently unequal.''
\end{quote}
\noindent --- Chief Justice Warren, for a unanimous Court, 347 U.S.\ at 495.

\begin{quote}
``To separate them from others of similar age and qualifications solely because of their race generates a feeling of inferiority as to their status in the community that may affect their hearts and minds in a way unlikely ever to be undone.''
\end{quote}
\noindent --- Chief Justice Warren, for a unanimous Court, 347 U.S.\ at 494.

\end{document}
